\newcommand{\equationcpointone}{\se(\none) \equiv []}
\newcommand{\equationcpointtwo}{\se(x \in \Y) \equiv x}
\newcommand{\equationcpointthree}{\se(\tup{a, b, \dots}) \equiv \se(a) \concat \se(b) \concat \dots}
\newcommand{\equationcpointfour}{\se(a, b, \dots) \equiv \se(\tup{a, b, \dots})}
\newcommand{\equationcpointfive}{
    \se_{l \in \N}\colon\left\{\begin{aligned}
        \N_{2^{8l}} \to \Y_l \\
        x \mapsto \begin{cases}
            [] \when l = 0 \\
            [x \bmod 256] \concat \se_{l - 1}\left(\left\lfloor \frac{x}{256} \right\rfloor\right) \otherwise
        \end{cases}
    \end{aligned}\right.
}
\newcommand{\equationcpointsix}{
    \se\colon\left\{\begin{aligned}
        \N_{2^{64}} \to \Y_{1:9} \\
        x \mapsto \begin{cases}
            [0] \when x = 0 \\
            \left[2^8-2^{8-l} + \ffrac{x}{2^{8l}}\right] \concat \se_l(x \bmod 2^{8l}) \when \exists l \in \N_8 : 2^{7l} \le x < 2^{7(l+1)} \\
            [2^8-1] \concat \se_8(x) \otherwhen x < 2^{64} \\
        \end{cases}
    \end{aligned}\right.
}
\newcommand{\equationcpointseven}{\se([i_0, i_1, ...]) \equiv \se(i_0) \concat \se(i_1) \concat \dots}
\newcommand{\equationcpointeight}{\var{x} \equiv \tup{|x|, x}\text{ thus }\se(\var{x}) \equiv \se(|x|)\concat\se(x)}
\newcommand{\equationcpointnine}{
    \begin{align}
        \maybe{x} \equiv \begin{cases}
            0 \when x = \none \\
            (1, x) \otherwise
        \end{cases}
    \end{align}
}
\newcommand{\equationcpointten}{
    \begin{align}
        \se(b \in \mathbb{B}) \equiv \begin{cases}
            [] \when b = [] \\
            \left[\sum\limits_{i=0}^{\min(8, |b|)}{b_i\cdot 2^i}\right] \concat \se(b_{8\dots}) \otherwise
        \end{cases}
    \end{align}
}
\newcommand{\equationcpointeleven}{\forall K, V: \se(d \in \dict{K}{V}) \equiv \se(\var{\orderby{k}{\tup{\se(k), \se(d[k])} \mid k \in \keys{d}}})}
\newcommand{\equationcpointtwelve}{\se(\{a, b, c, \dots\}) \equiv \se(a) \concat \se(b) \concat \se(c) \concat \dots \where a < b < c < \dots}
\newcommand{\equationcpointthirteen}{
    \begin{align}
        \se(\mathbf{B}) = \se\,\left\lparen
            \mathbf{H},\ \se_T(\xttickets),\ \se_P(\xtpreimages),\ \se_G(\xtguarantees),\ \se_A(\xtassurances),\ \se_D(\xtdisputes)
        \right\rparen
    \end{align}
}
\newcommand{\equationcpointfourteen}{
    \begin{align}
        \se_T(\xttickets) = \se(\var{\xttickets})
    \end{align}
}
\newcommand{\equationcpointfifteen}{
    \begin{align}
        \se_P(\xtpreimages) = \se(\var{[\tup{s, \var{p}} \mid \tup{s, p} \orderedin \xtpreimages]})
    \end{align}
}
\newcommand{\equationcpointsixteen}{
    \begin{align}
        \se_G(\xtguarantees) = \se(\var{[\tup{w, \se_4(t), \var{a}} \mid \tup{w, t, a} \orderedin \xtguarantees]})
    \end{align}
}
\newcommand{\equationcpointseventeen}{
    \begin{align}
        \se_A(\xtassurances) = \se(\var{[\tup{a, f, \se_2(v), s} \mid \tup{a, f, v, s} \orderedin \xtassurances]})
    \end{align}
}
\newcommand{\equationcpointeighteen}{
    \begin{align}
        \se_D((\mathbf{v}, \mathbf{c}, \mathbf{f})) = \se(\var{[(r, \se_4(a), [(v, \se_2(i), s) \mid (v, i, s) \orderedin \mathbf{j}]) \mid (r, a, \mathbf{j}) \orderedin \mathbf{v}]}, \var{\mathbf{c}}, \var{\mathbf{f}})
    \end{align}
}
\newcommand{\equationcpointnineteen}{
    \begin{align}
        \se(\mathbf{H}) = \se_U(\mathbf{H})\concat\se(\mathbf{H}_s)
    \end{align}
}
\newcommand{\equationcpointtwenty}{
    \begin{align}
        \se_U(\mathbf{H}) = \se(\mathbf{H}_p,\mathbf{H}_r,\mathbf{H}_x)\concat\se_4(\mathbf{H}_t)\concat\se(\maybe{\mathbf{H}_e},\maybe{\mathbf{H}_w},\var{\mathbf{H}_o},\se_2(\mathbf{H}_i),\mathbf{H}_v)
    \end{align}
}
\newcommand{\equationcpointtwentyone}{
    \begin{align}
        \se(x \in \mathbb{X}) \equiv \se(x_a, x_s, x_b, x_l)\concat\se_4(x_t)\concat\se(\var{x_\mathbf{p}})
    \end{align}
}
\newcommand{\equationcpointtwentytwo}{
    \begin{align}
        \se(x \in \mathbb{S}) \equiv \se(x_h) \concat \se_4(x_l)\concat\se(x_u, x_e) \concat \se_2(x_n)
    \end{align}
}
\newcommand{\equationcpointtwentythree}{
    \begin{align}
        \se(x \in \mathbb{L}) \equiv \se_4(x_s)\concat\se(x_c, x_l)\concat\se_8(x_g)\concat\se(O(x_\mathbf{o}))
    \end{align}
}
\newcommand{\equationcpointtwentyfour}{
    \begin{align}
        \se(x \in \mathbb{W}) \equiv \se(x_s, x_x, x_c, x_a, \var{x_\mathbf{o}}, \var{x_\mathbf{l}}, \var{x_\mathbf{r}})
    \end{align}
}
\newcommand{\equationcpointtwentyfive}{
    \begin{align}
        \se(x \in \mathbb{P}) \equiv \se(\var{x_\wp¬authtoken}, \se_4(x_\wp¬authcodehost), x_\wp¬authcodehash, \var{x_\wp¬authparam}, x_\wp¬context, \var{x_\wp¬workitems})
    \end{align}
}
\newcommand{\equationcpointtwentysix}{
    \begin{align}
        \se(x \in \mathbb{I}) \equiv \se(
            \se_4(x_s),
            x_c,
            \var{x_\mathbf{y}},
            \se_8(x_g),
            \se_8(x_a),
            \var{\se_I^\#(x_\mathbf{i})},
            \var{\sq{(h, \se_4(i)) \mid (h, i) \orderedin x_\mathbf{x}}},
            \se_2(x_e)
        )
    \end{align}
}
\newcommand{\equationcpointtwentyseven}{
    \begin{align}
        \se(x \in \mathbb{C}) \equiv \se(x_\mathbf{y}, x_r)
    \end{align}
}
\newcommand{\equationcpointtwentyeight}{
    \begin{align}
        \se(x \in \defxfer) \equiv \se(\se_4(x_s), \se_4(x_d), \se_8(x_a), \se(x_m), \se_8(x_g))
    \end{align}
}
\newcommand{\equationcpointtwentynine}{
    \begin{align}
        O(o \in \mathbb{J} \cup \Y) \equiv \begin{cases}
            (0, \var{o}) \when o \in \Y \\
            1 \when o = \infty \\
            2 \when o = \panic \\
            3 \when o = \badexports \\
            4 \when o = \token{BAD} \\
            5 \when o = \token{BIG} \\
        \end{cases}
    \end{align}
}
\newcommand{\equationcpointthirty}{
    \begin{align}
        \se_I(h \in \H \cup \H^\boxplus, i \in \N_{2^{15}}) \equiv \begin{cases}
            (h, \se_2(i)) \when h \in \H\\
            (r, \se_2(i + 2^{15})) \when \exists r \in \H, h = r^\boxplus
        \end{cases}
    \end{align}
}
